%% ---------------------------
% Zusätzliche Pakete und Definitionen
%% ---------------------------

\usepackage[utf8]{inputenc}
\usepackage[english]{babel} % Sprachanpassungen und Trennmuster
\usepackage[T1]{fontenc}    % T1 Schrift Encoding
\usepackage{textcomp}       % zusätzliche Symbole (Text Companion)
\usepackage{mathptmx}              %% --- Times mit Matheschriften
% \usepackage{mathpazo}            %% --- Palatino
\usepackage[scaled=.90]{helvet}    %% --- Helvetica

\usepackage{fontspec} % Font selection for XeLaTeX; see fontspec.pdf for documentation
\defaultfontfeatures{Mapping=tex-text} % to support TeX conventions like ``---''
\usepackage{xunicode} % Unicode support for LaTeX character names (accents, European chars, etc)
\usepackage{xltxtra} % Extra customizations for XeLaTeX

\setmainfont{Myriad Pro} % set the main body font (\textrm), assumes Charis SIL is installed
\setsansfont{Myriad Pro}

\usepackage{booktabs}
\usepackage{multimedia}
\usepackage{listings}


\usepackage{color}

\definecolor{pblue}{rgb}{0.13,0.13,1}
\definecolor{pgreen}{rgb}{0,0.5,0}
\definecolor{pred}{rgb}{0.9,0,0}
\definecolor{pgrey}{rgb}{0.46,0.45,0.48}

\lstset{language=Java,
  showspaces=false,
  showtabs=false,
  breaklines=true,
  showstringspaces=false,
  breakatwhitespace=true,
  commentstyle=\color{pgreen},
  keywordstyle=\color{pblue},
  stringstyle=\color{pred},
  basicstyle=\ttfamily\scriptsize,
  moredelim=[il][\textcolor{pgrey}]{$$},
  moredelim=[is][\textcolor{pgrey}]{\%\%}{\%\%}
}

\usepackage{enumerate}

%% ---------------------------
% Beamer Theme
%% ---------------------------
\usetheme{Esslingen}

%\usecolortheme{default}  %      kein Farb-Hintergrund in Blöcken
%\usecolortheme{lily}     %      kein Farb-Hintergrund in Blöcken
 \usecolortheme{rose}     %  dezenter Farb-Hintergrund in Blöcken
%\usecolortheme{orchid}   % kräftiger Farb-Hintergrund in Blöcken

\setbeamercovered{transparent=30}

%% ---------------------------
%% Definitionen für die Titelseite
%% ---------------------------
%% Die optionalen Argumente werden für die Fußzeile benutzt

\title[Software Architectures]{Software Architectures}

\subtitle[Design Patterns in Software Engineering]{Design Patterns in Software Engineering}

\author[Prof. Dr. Jörg Friedrich]{Prof. Dr. Jörg Friedrich\inst{1}}% \and
 %Zweiter Autor\inst{1} \and
 %Dritter Autor\inst{2}}

\titlegraphic{\includegraphics[scale=0.5]{figures/HEGoeppingen}}

\institute[University of Applied Sciences Esslingen]{%
\inst{1}University of Applied Sciences Esslingen, Faculty of Information Technology
%  \and
% \inst{2}Andere Einrichtung
}

\date[\copyright{} 2015]{\today}

